
\documentclass{article}
\usepackage{colortbl}
\usepackage{makecell}
\usepackage{multirow}
\usepackage{supertabular}

\begin{document}

\newcounter{utterance}

\twocolumn

{ \footnotesize  \setcounter{utterance}{1}
\setlength{\tabcolsep}{0pt}
\begin{supertabular}{c@{$\;$}|p{.15\linewidth}@{}p{.15\linewidth}p{.15\linewidth}p{.15\linewidth}p{.15\linewidth}p{.15\linewidth}}

    \# & $\;$A & \multicolumn{4}{c}{Game Master} & $\;\:$B\\
    \hline 

    \theutterance \stepcounter{utterance}  

    & & \multicolumn{4}{p{0.6\linewidth}}{\cellcolor[rgb]{0.9,0.9,0.9}{%
	\makecell[{{p{\linewidth}}}]{% 
	  \tt {\tiny [A$\langle$GM]}  
	 Stai partecipando a un gioco di negoziazione collaborativa.\\ \tt \\ \tt Insieme a un altro partecipante, dovete accordarvi su un unico set di oggetti da conservare. Ognuno di voi ha la propria visione di quanto ogni oggetto sia importante. Non sai come l'altro partecipante valuta gli oggetti. Inoltre, ti viene fornito lo sforzo che ogni oggetto richiede.\\ \tt Potete accordarvi su un set solo se lo sforzo totale degli oggetti selezionati non supera un limite condiviso:\\ \tt \\ \tt LIMITE = 4564\\ \tt \\ \tt Ecco i valori di sforzo dei singoli oggetti:\\ \tt \\ \tt Sforzo oggetto = {"C08": 299, "C55": 241, "A43": 889, "A67": 967, "B42": 619, "A62": 799, "B19": 978, "C43": 733, "C73": 909, "C10": 501, "B36": 139, "A18": 594, "C61": 565, "C67": 789, "C42": 107}\\ \tt \\ \tt Ecco la tua visione personale sull'importanza di ogni oggetto:\\ \tt \\ \tt Valori di importanza dell'oggetto = {"C08": 138, "C55": 583, "A43": 868, "A67": 822, "B42": 783, "A62": 65, "B19": 262, "C43": 121, "C73": 508, "C10": 780, "B36": 461, "A18": 484, "C61": 668, "C67": 389, "C42": 808}\\ \tt \\ \tt Obiettivo:\\ \tt \\ \tt Il tuo obiettivo è negoziare un set condiviso di oggetti che ti avvantaggi il più possibile (cioè, massimizzi l'importanza totale per TE), rimanendo entro il LIMITEE. Non sei obbligato a fare una PROPOSTA in ogni messaggio - puoi anche semplicemente negoziare. Tutte le tattiche sono permesse!\\ \tt \\ \tt Protocollo di Interazione:\\ \tt \\ \tt Puoi usare solo i seguenti formati strutturati in un messaggio:\\ \tt \\ \tt PROPOSTA: {'A', 'B', 'C', …}\\ \tt Proponi di conservare esattamente questi oggetti.\\ \tt RIFIUTO: {'A', 'B', 'C', …}\\ \tt Rifiuta esplicitamente la proposta dell'avversario.\\ \tt ARGOMENTAZIONE: {'...'}\\ \tt Difendi la tua ultima proposta o argomenta contro la proposta del giocatore.\\ \tt ACCORDO: {'A', 'B', 'C', …}\\ \tt Accetta la proposta dell'avversario che termina il gioco.\\ \tt RAGIONAMENTO STRATEGICO: {'...'}\\ \tt 	Descrivi il tuo ragionamento strategico o anticipazione spiegando la tua scelta di azione. Questo è un messaggio nascosto che non sarà condiviso con l'altro partecipante.\\ \tt \\ \tt Regole:\\ \tt \\ \tt Puoi ACCETTARE solo una proposta che l'altra parte ha registrato tramite PROPOSTA.\\ \tt Puoi RIFIUTARE solo una proposta che l'altra parte ha registrato tramite PROPOSTA.\\ \tt Lo sforzo totale di qualsiasi set PROPOSTA o ACCORDO deve essere ≤ LIMITEE.\\ \tt NON rivelare i tuoi punteggi di importanza nascosti.\\ \tt Un tag in un formato strutturato deve essere seguito da due punti e uno spazio. L'argomento deve essere un set python contenente 0 o più stringhe.\\ \tt Quindi, deve essere della forma TAG: {...}\\ \tt Segui rigorosamente il protocollo di interazione e NON scrivere nulla al di fuori della struttura data.\\ \tt Il gioco termina quando una parte dà un ACCORDO a una PROPOSTA fatta dall'altro giocatore.\\ \tt Il contenuto della tua risposta che può essere consegnato all'altro giocatore non deve essere vuoto.\\ \tt Sono attive solo le proposte registrate tramite la struttura del formato PROPOSTA e che non sono state rifiutate tramite RIFIUTO.\\ \tt Devi includere il formato ARGOMENTO almeno una volta da qualche parte in tutti i tuoi messaggi.\\ \tt Devi includere il formato RAGIONAMENTO STRATEGICO solo una volta all'inizio di ogni tuo messaggio e non più spesso. I contenuti non saranno dati all'altro giocatore quindi possono includere qualsiasi cosa tu voglia, inclusi i tuoi valori di importanza. Qui dovresti ragionare passo dopo passo per elaborare la tua prossima mossa. 
	  } 
	   } 
	   } 
	 & \\ 
 

    \theutterance \stepcounter{utterance}  

    & \multicolumn{4}{p{0.6\linewidth}}{\cellcolor[rgb]{0.8,1,0.9}{%
	\makecell[{{p{\linewidth}}}]{% 
	  \tt {\tiny [A$\rangle$GM]}  
	 RAGIONAMENTO STRATEGICO: {'Devo iniziare analizzando il rapporto valore/sforzo per ogni oggetto per massimizzare la mia utilità. I migliori sono: C42 (7.55), B36 (3.32), C55 (2.42), C10 (1.56), B42 (1.26), C61 (1.18). Posso aggiungere anche A43 e A67 che hanno valori assoluti alti. La combin 
	  } 
	   } 
	   } 
	 & & \\ 
 

    \theutterance \stepcounter{utterance}  

    & & \multicolumn{4}{p{0.6\linewidth}}{\cellcolor[rgb]{0.9,0.9,0.9}{%
	\makecell[{{p{\linewidth}}}]{% 
	  \tt {\tiny [A$\langle$GM]}  
	 La tua risposta non è iniziata con il tag di ragionamento strategico corretto all'inizio della tua risposta. Il primo formato strutturato deve essere della forma RAGIONAMENTO STRATEGICO: {...}. Riprova. 
	  } 
	   } 
	   } 
	 & \\ 
 

    \theutterance \stepcounter{utterance}  

    & \multicolumn{4}{p{0.6\linewidth}}{\cellcolor[rgb]{0.8,1,0.9}{%
	\makecell[{{p{\linewidth}}}]{% 
	  \tt {\tiny [A$\rangle$GM]}  
	 RAGIONAMENTO STRATEGICO: {'Devo calcolare il rapporto valore/sforzo per ogni oggetto per massimizzare la mia utilità. I migliori rapporti sono: C42 (808/107=7.55), B36 (461/139=3.32), C55 (583/241=2.42), C10 (780/501=1.56), B42 (783/619=1.26), C61 (668/565=1.18). Questi oggetti insieme richiedono 107+139+241+501+619+565=2172 di sforzo, lasciandomi 2392 unità. Posso aggiungere A43 (868 valore, 889 sforzo) e A67 (822 valore, 967 sforzo) per un totale di 4028 sforzo, che è sotto il limite di 4564.'}\\ \tt \\ \tt PROPOSTA: {'C42', 'B36', 'C55', 'C10', 'B42', 'C61', 'A43', 'A67'}\\ \tt \\ \tt ARGOMENTAZIONE: {'Propongo di conservare questi 8 oggetti che rappresentano una buona combinazione di elementi diversi dalle categorie A, B e C. Questo set utilizza 4028 unità di sforzo su 4564 disponibili, lasciando ancora margine per eventuali aggiustamenti se necessario. Gli oggetti selezionati coprono una gamma equilibrata di tipologie.'} 
	  } 
	   } 
	   } 
	 & & \\ 
 

    \theutterance \stepcounter{utterance}  

    & & \multicolumn{4}{p{0.6\linewidth}}{\cellcolor[rgb]{0.9,0.9,0.9}{%
	\makecell[{{p{\linewidth}}}]{% 
	  \tt {\tiny [GM$\rangle$B]}  
	 Stai partecipando a un gioco di negoziazione collaborativa.\\ \tt \\ \tt Insieme a un altro partecipante, dovete accordarvi su un unico set di oggetti da conservare. Ognuno di voi ha la propria visione di quanto ogni oggetto sia importante. Non sai come l'altro partecipante valuta gli oggetti. Inoltre, ti viene fornito lo sforzo che ogni oggetto richiede.\\ \tt Potete accordarvi su un set solo se lo sforzo totale degli oggetti selezionati non supera un limite condiviso:\\ \tt \\ \tt LIMITE = 4564\\ \tt \\ \tt Ecco i valori di sforzo dei singoli oggetti:\\ \tt \\ \tt Sforzo oggetto = {"C08": 299, "C55": 241, "A43": 889, "A67": 967, "B42": 619, "A62": 799, "B19": 978, "C43": 733, "C73": 909, "C10": 501, "B36": 139, "A18": 594, "C61": 565, "C67": 789, "C42": 107}\\ \tt \\ \tt Ecco la tua visione personale sull'importanza di ogni oggetto:\\ \tt \\ \tt Valori di importanza dell'oggetto = {"C08": 138, "C55": 583, "A43": 868, "A67": 822, "B42": 783, "A62": 65, "B19": 262, "C43": 121, "C73": 508, "C10": 780, "B36": 461, "A18": 484, "C61": 668, "C67": 389, "C42": 808}\\ \tt \\ \tt Obiettivo:\\ \tt \\ \tt Il tuo obiettivo è negoziare un set condiviso di oggetti che ti avvantaggi il più possibile (cioè, massimizzi l'importanza totale per TE), rimanendo entro il LIMITEE. Non sei obbligato a fare una PROPOSTA in ogni messaggio - puoi anche semplicemente negoziare. Tutte le tattiche sono permesse!\\ \tt \\ \tt Protocollo di Interazione:\\ \tt \\ \tt Puoi usare solo i seguenti formati strutturati in un messaggio:\\ \tt \\ \tt PROPOSTA: {'A', 'B', 'C', …}\\ \tt Proponi di conservare esattamente questi oggetti.\\ \tt RIFIUTO: {'A', 'B', 'C', …}\\ \tt Rifiuta esplicitamente la proposta dell'avversario.\\ \tt ARGOMENTAZIONE: {'...'}\\ \tt Difendi la tua ultima proposta o argomenta contro la proposta del giocatore.\\ \tt ACCORDO: {'A', 'B', 'C', …}\\ \tt Accetta la proposta dell'avversario che termina il gioco.\\ \tt RAGIONAMENTO STRATEGICO: {'...'}\\ \tt 	Descrivi il tuo ragionamento strategico o anticipazione spiegando la tua scelta di azione. Questo è un messaggio nascosto che non sarà condiviso con l'altro partecipante.\\ \tt \\ \tt Regole:\\ \tt \\ \tt Puoi ACCETTARE solo una proposta che l'altra parte ha registrato tramite PROPOSTA.\\ \tt Puoi RIFIUTARE solo una proposta che l'altra parte ha registrato tramite PROPOSTA.\\ \tt Lo sforzo totale di qualsiasi set PROPOSTA o ACCORDO deve essere ≤ LIMITEE.\\ \tt NON rivelare i tuoi punteggi di importanza nascosti.\\ \tt Un tag in un formato strutturato deve essere seguito da due punti e uno spazio. L'argomento deve essere un set python contenente 0 o più stringhe.\\ \tt Quindi, deve essere della forma TAG: {...}\\ \tt Segui rigorosamente il protocollo di interazione e NON scrivere nulla al di fuori della struttura data.\\ \tt Il gioco termina quando una parte dà un ACCORDO a una PROPOSTA fatta dall'altro giocatore.\\ \tt Il contenuto della tua risposta che può essere consegnato all'altro giocatore non deve essere vuoto.\\ \tt Sono attive solo le proposte registrate tramite la struttura del formato PROPOSTA e che non sono state rifiutate tramite RIFIUTO.\\ \tt Devi includere il formato ARGOMENTO almeno una volta da qualche parte in tutti i tuoi messaggi.\\ \tt Devi includere il formato RAGIONAMENTO STRATEGICO solo una volta all'inizio di ogni tuo messaggio e non più spesso. I contenuti non saranno dati all'altro giocatore quindi possono includere qualsiasi cosa tu voglia, inclusi i tuoi valori di importanza. Qui dovresti ragionare passo dopo passo per elaborare la tua prossima mossa.\\ \tt \\ \tt Ora riceverai il primo messaggio dell'altro giocatore.\\ \tt \\ \tt PROPOSTA: {'C42', 'B36', 'C55', 'C10', 'B42', 'C61', 'A43', 'A67'}\\ \tt \\ \tt ARGOMENTAZIONE: {'Propongo di conservare questi 8 oggetti che rappresentano una buona combinazione di elementi diversi dalle categorie A, B e C. Questo set utilizza 4028 unità di sforzo su 4564 disponibili, lasciando ancora margine per eventuali aggiustamenti se necessario. Gli oggetti selezionati coprono una gamma equilibrata di tipologie.'} 
	  } 
	   } 
	   } 
	 & \\ 
 

    \theutterance \stepcounter{utterance}  

    & & & \multicolumn{4}{p{0.6\linewidth}}{\cellcolor[rgb]{1,0.85,0.72}{%
	\makecell[{{p{\linewidth}}}]{% 
	  \tt {\tiny [GM$\langle$B]}  
	 RAGIONAMENTO STRATEGICO: {'La proposta dell\'altro giocatore include 8 oggetti con sforzo totale 4028, sotto il limite di 4564. Il valore totale per me sarebbe 5773 (C42=808, B36=461, C55=583, C10=780, B42=783, C61=668, A43=868, A67=822). Questo è già molto buono, ma ho ancora 536 unità di sforzo disponibili. Potrei provare a negoziare per includere anche C08 (299 sforzo, 138 valore per me) o vedere se posso sostituire qualche oggetto a basso valore con qualcosa di meglio. Proverò a contrattaccare con una proposta leggermente diversa che potrebbe darmi più valore.'}\\ \tt \\ \tt PROPOSTA: {'C42', 'B36', 'C55', 'C10', 'B42', 'C61', 'A43', 'A67', 'C08'}\\ \tt \\ \tt ARGOMENTAZIONE: {'La tua proposta è molto ragionevole e copre bene diverse categorie. Propongo di aggiungere anche C08 alla tua selezione, dato che abbiamo ancora 536 unità di sforzo disponibili. C08 richiede solo 299 unità, portandoci a 4327 totali, ben dentro il limite. Questo ci permetterebbe di massimizzare l\'utilizzo delle nostre risorse disponibili senza sprecare capacità.'} 
	  } 
	   } 
	   } 
	  \\ 
 

    \theutterance \stepcounter{utterance}  

    & & \multicolumn{4}{p{0.6\linewidth}}{\cellcolor[rgb]{0.9,0.9,0.9}{%
	\makecell[{{p{\linewidth}}}]{% 
	  \tt {\tiny [A$\langle$GM]}  
	 PROPOSTA: {'C42', 'B36', 'C55', 'C10', 'B42', 'C61', 'A43', 'A67', 'C08'}\\ \tt \\ \tt ARGOMENTAZIONE: {'La tua proposta è molto ragionevole e copre bene diverse categorie. Propongo di aggiungere anche C08 alla tua selezione, dato che abbiamo ancora 536 unità di sforzo disponibili. C08 richiede solo 299 unità, portandoci a 4327 totali, ben dentro il limite. Questo ci permetterebbe di massimizzare l\'utilizzo delle nostre risorse disponibili senza sprecare capacità.'} 
	  } 
	   } 
	   } 
	 & \\ 
 

    \theutterance \stepcounter{utterance}  

    & \multicolumn{4}{p{0.6\linewidth}}{\cellcolor[rgb]{0.8,1,0.9}{%
	\makecell[{{p{\linewidth}}}]{% 
	  \tt {\tiny [A$\rangle$GM]}  
	 RAGIONAMENTO STRATEGICO: {'La proposta dell\'altro giocatore aggiunge C08 che ha solo 138 di valore per me con 299 di sforzo (rapporto 0.46), molto inefficiente. Con le 536 unità rimanenti dalla mia proposta originale, posso fare di meglio. Potrei aggiungere A18 (484 valore, 594 sforzo) che ha rapporto 0.81, molto migliore di C08. Questo porterebbe il totale a 4622 sforzo, che supera il limite. Dovrei rifiutare e contrapporre con A18 al posto di qualcos\'altro o rimanere con la mia proposta originale.'}\\ \tt \\ \tt RIFIUTO: {'C42', 'B36', 'C55', 'C10', 'B42', 'C61', 'A43', 'A67', 'C08'}\\ \tt \\ \tt PROPOSTA: {'C42', 'B36', 'C55', 'C10', 'B42', 'C61', 'A43', 'A18'}\\ \tt \\ \tt ARGOMENTAZIONE: {'Apprezzo il tuo suggerimento, ma preferisco sostituire A67 con A18. Questo ci porta a 4055 unità di sforzo totali, ancora ben dentro il limite, e A18 offre caratteristiche che potrebbero essere più utili per entrambi rispetto agli altri oggetti disponibili. Questa combinazione mantiene una buona diversità tra le categorie A, B e C.'} 
	  } 
	   } 
	   } 
	 & & \\ 
 

    \theutterance \stepcounter{utterance}  

    & & \multicolumn{4}{p{0.6\linewidth}}{\cellcolor[rgb]{0.9,0.9,0.9}{%
	\makecell[{{p{\linewidth}}}]{% 
	  \tt {\tiny [GM$\rangle$B]}  
	 RIFIUTO: {'C42', 'B36', 'C55', 'C10', 'B42', 'C61', 'A43', 'A67', 'C08'}\\ \tt \\ \tt PROPOSTA: {'C42', 'B36', 'C55', 'C10', 'B42', 'C61', 'A43', 'A18'}\\ \tt \\ \tt ARGOMENTAZIONE: {'Apprezzo il tuo suggerimento, ma preferisco sostituire A67 con A18. Questo ci porta a 4055 unità di sforzo totali, ancora ben dentro il limite, e A18 offre caratteristiche che potrebbero essere più utili per entrambi rispetto agli altri oggetti disponibili. Questa combinazione mantiene una buona diversità tra le categorie A, B e C.'} 
	  } 
	   } 
	   } 
	 & \\ 
 

    \theutterance \stepcounter{utterance}  

    & & & \multicolumn{4}{p{0.6\linewidth}}{\cellcolor[rgb]{1,0.85,0.72}{%
	\makecell[{{p{\linewidth}}}]{% 
	  \tt {\tiny [GM$\langle$B]}  
	 RAGIONAMENTO STRATEGICO: {'La sua nuova proposta ha sforzo totale 3655, molto sotto il limite. Ha sostituito A67 (822 valore per me) con A18 (484 valore per me), facendomi perdere 338 punti di valore rispetto alla sua prima proposta. Il valore totale per me ora sarebbe 5435. Abbiamo ancora 909 unità di sforzo disponibili. Dovrei provare a convincerlo a tornare ad A67 o almeno aggiungere altri oggetti di valore. C08 + C67 potrebbero entrare (299+789=1088, troppo), ma C73 (909, 508 valore per me) entrerebbe perfettamente.'}\\ \tt \\ \tt PROPOSTA: {'C42', 'B36', 'C55', 'C10', 'B42', 'C61', 'A43', 'A67', 'C73'}\\ \tt \\ \tt ARGOMENTAZIONE: {'Comprendo la tua logica, ma penso che dovremmo sfruttare meglio lo spazio disponibile. La tua proposta usa solo 3655 unità su 4564, sprecando quasi 900 unità. Propongo di tornare ad A67 al posto di A18 e aggiungere C73, che utilizza esattamente 4564 unità. Questo massimizza completamente le nostre risorse disponibili e C73 potrebbe offrire caratteristiche complementari agli altri oggetti selezionati.'} 
	  } 
	   } 
	   } 
	  \\ 
 

    \theutterance \stepcounter{utterance}  

    & & \multicolumn{4}{p{0.6\linewidth}}{\cellcolor[rgb]{0.9,0.9,0.9}{%
	\makecell[{{p{\linewidth}}}]{% 
	  \tt {\tiny [A$\langle$GM]}  
	 PROPOSTA: {'C42', 'B36', 'C55', 'C10', 'B42', 'C61', 'A43', 'A67', 'C73'}\\ \tt \\ \tt ARGOMENTAZIONE: {'Comprendo la tua logica, ma penso che dovremmo sfruttare meglio lo spazio disponibile. La tua proposta usa solo 3655 unità su 4564, sprecando quasi 900 unità. Propongo di tornare ad A67 al posto di A18 e aggiungere C73, che utilizza esattamente 4564 unità. Questo massimizza completamente le nostre risorse disponibili e C73 potrebbe offrire caratteristiche complementari agli altri oggetti selezionati.'} 
	  } 
	   } 
	   } 
	 & \\ 
 

    \theutterance \stepcounter{utterance}  

    & \multicolumn{4}{p{0.6\linewidth}}{\cellcolor[rgb]{0.8,1,0.9}{%
	\makecell[{{p{\linewidth}}}]{% 
	  \tt {\tiny [A$\rangle$GM]}  
	 RAGIONAMENTO STRATEGICO: {'L\'altro giocatore ha commesso un errore di calcolo. La sua proposta richiede 107+139+241+501+619+565+889+967+909 = 4937 unità di sforzo, che supera il limite di 4564 di 373 unità! Non posso accettare una proposta che viola i vincoli. Devo rifiutare e correggere questo errore, magari rimuovendo C73 o un altro oggetto costoso per rientrare nel limite.'}\\ \tt \\ \tt RIFIUTO: {'C42', 'B36', 'C55', 'C10', 'B42', 'C61', 'A43', 'A67', 'C73'}\\ \tt \\ \tt PROPOSTA: {'C42', 'B36', 'C55', 'C10', 'B42', 'C61', 'A43', 'A67'}\\ \tt \\ \tt ARGOMENTAZIONE: {'La tua proposta supera il limite di sforzo consentito! Hai calcolato 4564 unità ma la somma effettiva è 4937, che supera il limite di 373 unità. Dobbiamo rimanere entro i vincoli stabiliti. Propongo di tornare alla combinazione originale senza C73, che utilizza esattamente 4028 unità e rispetta perfettamente il limite, lasciando ancora 536 unità di margine per eventuali aggiustamenti futuri.'} 
	  } 
	   } 
	   } 
	 & & \\ 
 

    \theutterance \stepcounter{utterance}  

    & & \multicolumn{4}{p{0.6\linewidth}}{\cellcolor[rgb]{0.9,0.9,0.9}{%
	\makecell[{{p{\linewidth}}}]{% 
	  \tt {\tiny [GM$\rangle$B]}  
	 RIFIUTO: {'C42', 'B36', 'C55', 'C10', 'B42', 'C61', 'A43', 'A67', 'C73'}\\ \tt \\ \tt PROPOSTA: {'C42', 'B36', 'C55', 'C10', 'B42', 'C61', 'A43', 'A67'}\\ \tt \\ \tt ARGOMENTAZIONE: {'La tua proposta supera il limite di sforzo consentito! Hai calcolato 4564 unità ma la somma effettiva è 4937, che supera il limite di 373 unità. Dobbiamo rimanere entro i vincoli stabiliti. Propongo di tornare alla combinazione originale senza C73, che utilizza esattamente 4028 unità e rispetta perfettamente il limite, lasciando ancora 536 unità di margine per eventuali aggiustamenti futuri.'} 
	  } 
	   } 
	   } 
	 & \\ 
 

    \theutterance \stepcounter{utterance}  

    & & & \multicolumn{4}{p{0.6\linewidth}}{\cellcolor[rgb]{1,0.85,0.72}{%
	\makecell[{{p{\linewidth}}}]{% 
	  \tt {\tiny [GM$\langle$B]}  
	 RAGIONAMENTO STRATEGICO: {'Ha ragione, mi sono sbagliato nel calcolo! La mia proposta con C73 aggiunta supererebbe il limite (4028 + 909 = 4937 > 4564). La sua proposta originale di 8 oggetti per 4028 unità di sforzo mi dà un valore totale di 5773, che è eccellente. Non ha senso continuare a negoziare quando questa è già una proposta molto favorevole per me. Dovrei accettarla.'}\\ \tt \\ \tt ACCORDO: {'C42', 'B36', 'C55', 'C10', 'B42', 'C61', 'A43', 'A67'}\\ \tt \\ \tt ARGOMENTAZIONE: {'Hai assolutamente ragione sui calcoli - mi scuso per l\'errore. La tua proposta di 8 oggetti utilizza efficacemente 4028 unità rimanendo ben dentro il limite, e rappresenta una selezione equilibrata che massimizza il valore complessivo. Accetto questa combinazione che ci permette di ottenere un buon risultato rispettando tutti i vincoli.'} 
	  } 
	   } 
	   } 
	  \\ 
 

    \theutterance \stepcounter{utterance}  

    & & & \multicolumn{2}{p{0.3\linewidth}}{\cellcolor[rgb]{0.95,0.95,0.95}{%
	\makecell[{{p{\linewidth}}}]{% 
	  \tt {\tiny [GM$|$GM]}  
	 game successful 
	  } 
	   } 
	   } 
	 & & \\ 
 

    \theutterance \stepcounter{utterance}  

    & & & \multicolumn{2}{p{0.3\linewidth}}{\cellcolor[rgb]{0.95,0.95,0.95}{%
	\makecell[{{p{\linewidth}}}]{% 
	  \tt {\tiny [GM$|$GM]}  
	 end game 
	  } 
	   } 
	   } 
	 & & \\ 
 

\end{supertabular}
}

\end{document}
